%قبل از شروع مسئله ترجیحاً چیزی ننویسید. اما اگر خواستید بنویسید، در قالب لیستی (فقرات و شمارش) بنویسید.
%برای اضافه کردن نمره، بعدی از شروع‌مسئله، [x] را اضافه کنید و به جای x  نمره مورد نظر را قرار دهید. برای زیرمسئله‌ها هم به همین صورت است.
%\شروع‌مسئله[10]
\شروع‌مسئله [12]
%\متن‌سیاه{سوال عملی اول:}\\
%\متن‌سیاه{دراپ‌تست}\\
محمد در یک مسابقه‌ی معلومات علمی شرکت کرده است. مسابقه شامل \متن‌لاتین{n} سوال است. محمد سوال \متن‌لاتین{i} ام را به احتمال \متن‌لاتین{p[i]} درصد درست پاسخ می‌دهد و به احتمال \متن‌لاتین{100 - p[i]} درصد غلط پاسخ می‌دهد. هم چنین اگر محمد سوال \متن‌لاتین{i} ام را درست پاسخ دهد، به اندازه‌ی \متن‌لاتین{c[i]} تومان جایزه می‌گیرد. حال روند مسابقه بدین شکل است که سوال ها به ترتیب از محمد پرسیده می شود. هر بار اگر محمد سوال را درست پاسخ دهد، جایزه آن سوال را می‌گیرد و در صورت وجود، به سوال بعد می‌رود. اما اگر یک سوال را غلط پاسخ دهد، مسابقه تمام می شود و محمد جایزه هایی که تا کنون جمع کرده است را به خانه می‌برد. حال ما می‌توانیم برای کمک به محمد، ترتیب سوال ها را تغییر دهیم. هدف این است که پس از تغییر ترتیب سوال ها، امیدریاضی جایزه‌ای که محمد می‌برد ماکسیمم شود. حال محمد از شما پرسیده است که به ازای همه‌ی ترتیب های ممکن، ماکسیمم امیدریاضی جایزه چقدر است ؟

\متن‌سیاه{ورودی:}\\
ابتدا در یک خط تعداد سوالات یعنی \متن‌لاتین{n} می‌آید.\\
سپس \متن‌لاتین{n} خط می‌آید که هر کدام شامل دو عدد  \متن‌لاتین{c[i]} و  \متن‌لاتین{p[i]} است.\\
\متن‌سیاه{خروجی:}\\
در تنها سطر خروجی، عدد خواسته شده را با دقت ۳ رقم اعشار چاپ نمایید.\\
\متن‌سیاه{محدودیت‌ها:}\\
$1 \leq n \leq 10^{5}$\\
$0 \leq p[i] \leq 100$\\
$0 \leq c[i] \leq 10^{4}$\\

\متن‌سیاه{ورودی و خروجی نمونه:}\\

تست اول
\begin{latin}
	\begin{tabularx}{\textwidth}{|X | X|}
		\hline
		\textbf{stdin} &
		\textbf{stdout}\\
		\hline
		4 & 2.434\\
		100 1&\\
		80 1 &\\
		66 1&\\
		20 1&\\
		\hline
	\end{tabularx}
\end{latin}

تست دوم
\begin{latin}
	\begin{tabularx}{\textwidth}{|X | X|}
		\hline
		\textbf{stdin} &
		\textbf{stdout}\\
		\hline
		3 & 2.326\\
		50 1&\\
		33 6&\\
		70 1&\\
		\hline
		
	\end{tabularx}
\end{latin}
%}
% صورت اصلی مسئله را در این‌جا بنویسید.
%اگر مسئله چند قسمت داشت به این صورت عمل کنید:
%\شروع{زیرمسائل}
%\زیرمسئله[0] % مسئلۀ ۱آ
%\زیرمسئله % مسئلۀ ۱ب
%\زیرمسئله % مسئلۀ ۱پ
%\زیرمسئله % مسئلۀ ۱ت
%\پایان{زیرمسائل}

\solution{
	پاسخ نمونه.
}




%%%%%%%%%%%%%%%%%%%%%%%%%%%%%%%%%%%%%%%%%%%%%%%%%%%%%
% برخی دستورات مفید برای تمرین‌ها. %
%%%%%%%%%%%%%%%%%%%%%%%%%%%%%%%%%%%%%%%%%%%%%%%%%%%%%

% فرمول‌های نمونه:
%$f(x) = \Theta(x)$
%$T(x, y) \leq \log(x) + 2^y + \binom{2n}{n}$
% {\tt code\_function}


% لیست موارد بدون شماره‌گذاری:
%\شروع{فقرات}
%    \فقره مورد اول 
%        \شروع{فقرات}
%            \فقره مورد اول داخلی 
%                \شروع{فقرات}
%                    \فقره اولین مورد خیلی داخلی 
%                    \فقره دومین مورد خیلی داخلی 
%                \پایان{فقرات}
%            \فقره مورد دوم داخلی 
%        \پایان{فقرات}
%    \فقره مورد دوم 
%\پایان{فقرات}

% لیست موارد شماره‌گذاری شده:
%\شروع{شمارش}
%    \فقره یک مورد 
%    \فقره یک مورد دیگر 
%    \فقره باز هم یک مورد دیگر
%\پایان{شمارش}

% برای نوشتن یک فرمول چند خطی که یک شماره بخوره. برای حالت‌های دیگر به سایت زیر مراجعه کنید:
% https://www.overleaf.com/learn/latex/Aligning_equations_with_amsmath
%\begin{align} 
%    \begin{split}
%        (x+y)^3 &= (x+y)^2(x+y) \\
%                &= (x^2+2xy+y^2)(x+y) \\
%                &= (x^3+2x^2y+xy^2) + (x^2y+2xy^2+y^3) \\
%                &= x^3+3x^2y+3xy^2+y^3
%    \end{split}                                 
%\end{align}

% نوشتن ماتریس و آرایه دوبعدی
%\begin{align}
%    A = 
%    \begin{bmatrix}
%        A_{11} & A_{21} \\
%        A_{21} & A_{22}
%    \end{bmatrix}
%\end{align}

% اضافه کردن تصویر با شرح. برای اضافه کردن تصویر بدون شرح، فقط دستور درج‌تصویر را استفاده کنید.
%\شروع{شکل}[h!]
%\تنظیم‌ازوسط
%  \درج‌تصویر[پهنا=0.5\پهنای‌متن]{img.jpg}
%\شرح{شرح تصویر}
%\برچسب{تصویر:۱}
%\پایان{شکل}
%\پایان{زیرمسائل}

%کد زیر را ببینید.
%%\lstinputlisting{source_filename.py}
%\begin{latin}
%\begin{lstlisting}
%# Import the modules
%import sys
%import random
%
%ans = True
%
%while ans:
%question = raw_input("Ask the magic 8 ball a question: (press enter to quit) ")
%
%answers = random.randint(1,8)
%
%if question == "":
%sys.exit()
%
%elif answers == 1:
%print "It is certain"
%
%elif answers == 2:
%print "Outlook good"
%
%elif answers == 3:
%print "You may rely on it"
%
%elif answers == 4:
%print "Ask again later"
%
%elif answers == 5:
%print "Concentrate and ask again"
%
%elif answers == 6:
%print "Reply hazy, try again"
%
%elif answers == 7:
%print "My reply is no"
%
%elif answers == 8:
%print "My sources say no"
%\end{lstlisting}
%\end{latin}